\documentclass[11pt]{article}
\usepackage[a4paper,margin=25mm]{geometry}
\usepackage{amsmath,amssymb,amsthm,booktabs,array}
\usepackage[T1]{fontenc}
\usepackage[colorlinks=true,linkcolor=blue,citecolor=blue,urlcolor=blue]{hyperref}
\newtheorem{theorem}{Theorem}
\newtheorem{lemma}{Lemma}
\newtheorem{corollary}{Corollary}
\newcommand{\Tr}{\operatorname{Tr}}
\newcommand{\diag}{\operatorname{diag}}
\newcommand{\ket}[1]{\lvert #1\rangle}
\newcommand{\bra}[1]{\langle #1\rvert}
\title{Four-level closure separation at Narasimhachar--Gour's Conjecture 1\newline
\large Supplementary Sec.\ S2 F, p.\ S11:\newline
beyond Maestro's exact qutrit separation (2026),\newline
using Dykema--Juschenko (2009)}
\author{Haobo Ma\thanks{ChronoAI Pte Ltd, Singapore.}\and
Wenlin Zhang\thanks{National University of Singapore; ChronoAI Pte Ltd, Singapore.}}
\date{4 November 2026}
\begin{document}
\maketitle
\begin{abstract}
Extending Maestro's exact qutrit separation (2026) by the unitary-moment
obstruction of Dykema and Juschenko (2009),
we refute Conjecture 1 of Narasimhachar and Gour (arXiv:1409.7740v3,
Supplementary Sec.\ S2 F, p.\ S11) even in the closure of finite
ground-space thermal transitions. Two diagonal Kraus operators send the uniform pure state of a
four-level system with energies $0,1,3,7$ to a rank-two state with the same
populations. Any energy-conserving realization with an ancilla initially
maximally mixed on its full ground eigenspace would express the output
correlation matrix as normalized traces of four unitaries. Two linear
relations make that representation impossible, with a dimension-independent
gap. The trace-norm error of any such thermal output exceeds $10^{-8}$.
A three-vector projection obstruction also refutes Conjecture 2 in
versions 1 and 2 of the posing source. The four-level obstruction extends
to every system with at least four occupied distinct energy levels.
Together with the known qutrit closure theorem,
this makes four the minimum dimension for closure separation under
distinct-level and distinct-gap assumptions.
\end{abstract}
\noindent\textbf{License.} Creative Commons Attribution 4.0 International (CC BY 4.0).

\section{The conjecture and the obstruction}
Narasimhachar and Gour's Conjecture 1 (Supplementary Sec.\ S2 F, p.\ S11)
concerns equivalence of cooling maps
and low-temperature thermal operations for the feasibility of \emph{state
transitions}, not just equality of channel sets \cite[Sec.\ S2 F,
Conjecture 1]{NG}. Exact finite-bath separation is already established
for qutrits by Maestro \cite[Letter, first main result; Supplement,
pp.\ 1--2]{MA}, while every generic qutrit cooling channel remains
approximable. We give a rational four-level cooling transition that no exact finite
ground-space thermal operation realizes. Throughout, this thermal model means
\begin{equation}\label{eq:thermal}
\mathcal T(\rho)=\Tr_A\!\left[U(\rho\otimes P_0/g)U^\dagger\right],
\qquad [U,H_S\otimes I+I\otimes H_A]=0,
\end{equation}
where $H_A$ acts on a nonzero finite-dimensional space, $P_0$ is its full
ground-space projector, and $g=\Tr P_0\geq1$. This is the idealized ground-state
limit in \cite[main-text Eq.\ (6); Supplementary Sec.\ S2 A]{NG}.
With $E_0<\cdots<E_3$, a cooling map has Kraus operators that are diagonal
or proportional to $\ket j\bra k$ with $j<k$.

The algebraic obstruction predates the thermodynamic conjecture.
Dykema and Juschenko \cite[Proposition 2.10 and Corollary 2.11]{DJ}
exhibit a rank-two $4\times4$ correlation matrix outside the set of
tracial unitary moments. Haagerup and Musat
\cite[Lemma 2.8, Example 3.2 and Remark 3.6]{HM} relate this obstruction to
nonfactorizable Schur multipliers. Their Remark 4.8(2) also records an
earlier four-level example of K\"ummerer (1986). Thus the recorded
chronology is K\"ummerer (1986), Dykema--Juschenko (2009), then
Haagerup--Musat (2010 preprint; 2011 journal).
We use a rational rank-two variant and
give the elementary matrix argument below. The thermodynamic step is that
unchanged, strictly positive populations force the four ground-space
blocks of any realization to be unitary. This turns a single state transition
into precisely the unitary-moment problem. We extend its elementary proof
to approximate thermal transitions and to the earlier vector conjecture
\cite[Supplementary Sec.\ II F, Conjecture 2, p.\ S11]{NG1,NG2}, also posed
by Narasimhachar in August 2014 \cite{MSE}.

Related closure separations at positive temperature are known for an equally
spaced qutrit \cite[Sec.\ IV C, Eq.\ (36)]{DHD} and for paired coherences in
a four-level system \cite[Sec.\ III B, Eqs.\ (20)--(28)]{CS}.
Those results use finite-temperature Gibbs baths, rather than
\eqref{eq:thermal}. Schur factorization also has earlier thermodynamic
applications \cite[Appendix C.2]{LSBNW}. The rational projection-trace
obstruction underlying Maestro's example appears in \cite[Proposition 3.2
and Example 3.4]{MR}; a broader two-unitary rational-weight criterion is
given in \cite[Proposition 3.6]{MUS}. These prior results are distinct from
the uniform four-unitary obstruction used below.

\section{An explicit transition}
Index the energy eigenbasis by $0,1,2,3$ and put
\begin{equation}\label{eq:data}
H_S=\diag(0,1,3,7),\qquad
B=\begin{pmatrix}1&0\\0&1\\3/5&4/5\\3/5&4i/5\end{pmatrix},
\qquad C=BB^\dagger.
\end{equation}
Let $\ket+=\frac12\sum_{j=0}^3\ket j$, $\rho=\ket+\bra+$, and $\sigma=C/4$.
The four rows of $B$ have norm one, so $C$ is positive semidefinite with
unit diagonal. Its rank is two, since its first two rows are independent.
Consequently $\sigma$ is a density matrix, with all four populations $1/4$.
Define
\begin{equation}\label{eq:kraus}
K_0=\diag(1,0,3/5,3/5),\qquad
K_1=\diag(0,1,4/5,4i/5).
\end{equation}
Then $K_0^\dagger K_0+K_1^\dagger K_1=I$ and
$\sum_{r=0}^1K_r\rho K_r^\dagger=C/4$ entry by entry. Thus these diagonal
Kraus operators define a cooling map realizing $\rho\mapsto\sigma$.
The positive energy gaps are $1,2,3,4,6,7$, all different, as required by
the source's nondegeneracy assumptions.
For this pair, the matrix $Q$ of the source's Theorem 1 (Supplementary
Theorem S2.4) is exactly $C$; its cooling-map criterion is satisfied.

\begin{lemma}\label{lem:blocks}
Let $H_S=\diag(E_0,\ldots,E_3)$ with strictly increasing energies. Let
$\ket\psi=\sum_k\alpha_k\ket k$ with $\alpha_k\ne0$ for every $k$.
If a map of the form \eqref{eq:thermal} preserves the populations of
$\ket\psi\bra\psi$, then its output has entries
\begin{equation}\label{eq:gram}
\mathcal T(\ket\psi\bra\psi)_{jk}
=\alpha_j\overline{\alpha_k}\,g^{-1}\Tr(V_jV_k^\dagger)
\end{equation}
for four $g\times g$ unitary matrices $V_0,\ldots,V_3$.
\end{lemma}
\begin{proof}
Choose an energy eigenbasis of the ancilla, with energies $e_a$ and ground
energy $e_0$, and write $G=\{a:e_a=e_0\}$. Put
$u_{jb,ka}=\bra{j,b}U\ket{k,a}$. Energy conservation gives
\begin{equation}\label{eq:support}
u_{jb,ka}=0\quad\text{unless}\quad E_j+e_b=E_k+e_a.
\end{equation}
For $a\in G$, a nonzero entry requires $j\leq k$.
For fixed $j,b$ and $a\in G$,
at most one $k$ satisfies \eqref{eq:support}, because the $E_k$ are distinct.
Hence the cross terms in each output population vanish. With
\[
P_{j|k}=\frac1g\sum_{a\in G}\sum_b|u_{jb,ka}|^2,
\qquad p_k=|\alpha_k|^2>0,
\]
the populations transform as $p\mapsto Pp$; $P$ is nonnegative,
column stochastic by unitarity of $U$, and $P_{j|k}=0$ for $j>k$.
Population preservation implies
\[
0=\sum_kE_kp_k-\sum_jE_j(Pp)_j
=\sum_{j<k}(E_k-E_j)P_{j|k}p_k.
\]
All terms are nonnegative and every coefficient of $P_{j|k}$ is positive,
so $P_{j|k}=0$ for $j<k$. Thus $P=I$. Each vanished sum of squares forces
$u_{jb,ka}=0$ for $a\in G$ and $j\ne k$. When $j=k$,
\eqref{eq:support} forces $b\in G$. Consequently
\[
U\ket{k,a}=\ket k\otimes V_k\ket a\qquad(a\in G).
\]
The $g$ columns on the right are orthonormal and lie in the $g$-dimensional
ground space; each $V_k$ is therefore unitary. Substituting these blocks
in \eqref{eq:thermal} gives \eqref{eq:gram}.
\end{proof}

\begin{theorem}\label{thm:counterexample}
The cooling transition \eqref{eq:data}--\eqref{eq:kraus} has no realization
of the form \eqref{eq:thermal}, for any finite ancilla Hamiltonian or ground
degeneracy.
\end{theorem}
\begin{proof}
Suppose it did. Lemma~\ref{lem:blocks} gives unitaries with
$g^{-1}\Tr(V_jV_k^\dagger)=C_{jk}$. In the Hilbert space of matrices with
inner product $\langle X,Y\rangle=g^{-1}\Tr(XY^\dagger)$, the four $V_j$
therefore have the same Gram matrix as the rows of $B$. Set $a=3/5$, $b=4/5$.
The row relations $B_2=aB_0+bB_1$ and $B_3=aB_0+ibB_1$ imply
\begin{equation}\label{eq:relations}
V_2=aV_0+bV_1,\qquad V_3=aV_0+ibV_1.
\end{equation}
Indeed, expanding the squared norm of either residual uses only the Gram
entries, so gives the same value, zero, as for the corresponding row residual
of $B$; positive definiteness of the matrix inner product makes the residual
itself zero. Put $X=V_0V_1^\dagger$. Since $a^2+b^2=1$, unitarity in
\eqref{eq:relations} gives
\[
I=V_2V_2^\dagger=I+ab(X+X^\dagger),\qquad
I=V_3V_3^\dagger=I+iab(X^\dagger-X).
\]
As $ab\ne0$, these equations yield $X+X^\dagger=0$ and
$X^\dagger-X=0$, hence $X=0$. But $X$ is unitary on a nonzero space,
a contradiction.
\end{proof}

\section{A uniform gap and approximate transitions}
Write $\|A\|_{2,g}=(\Tr(AA^\dagger)/g)^{1/2}$ for the normalized
Hilbert--Schmidt norm.
\begin{lemma}\label{lem:gap}
For any four $g\times g$ unitaries, $g\geq1$, their correlation matrix
$D_{jk}=g^{-1}\Tr(V_jV_k^\dagger)$ satisfies
$\max_{j,k}|D_{jk}-C_{jk}|\geq1/144$.
\end{lemma}
\begin{proof}
Put $\delta=\max|D_{jk}-C_{jk}|$, $a=3/5$, $b=4/5$.
The squared norms of the two residuals in \eqref{eq:relations} vanish
for $C$. Expanding them with $D$ therefore bounds each by
$(1+a+b)^2\delta$, so each residual norm is at most $12\sqrt\delta/5$.
For $A=aV_0+bV_1$ one has $\|A\|_\infty\leq7/5$ and thus
\[
\|AA^\dagger-I\|_{2,g}
\leq(\|A\|_\infty+1)\|A-V_2\|_{2,g}
\leq144\sqrt\delta/25.
\]
The same estimate holds for $aV_0+ibV_1$ and $V_3$.
Dividing the resulting unitarity identities by $ab=12/25$ gives
$\|X+X^\dagger\|_{2,g}\leq12\sqrt\delta$ and
$\|X^\dagger-X\|_{2,g}\leq12\sqrt\delta$, where $X=V_0V_1^\dagger$.
The triangle inequality gives $1=\|X\|_{2,g}\leq12\sqrt\delta$.
\end{proof}

\begin{corollary}\label{cor:robust}
For the data \eqref{eq:data}, every output $\tau=\mathcal T(\rho)$
of \eqref{eq:thermal} satisfies $\|\tau-\sigma\|_1>10^{-8}$.
In particular, $\sigma$ is outside the closure of these thermal outputs.
\end{corollary}
\begin{proof}
Let $P_{j|k}$ be as in Lemma~\ref{lem:blocks}, and set
$\ell_k=1-P_{k|k}$. The energy loss
$L=\Tr(H_S(\rho-\tau))$ satisfies
\[
L=\tfrac14\sum_{j<k}(E_k-E_j)P_{j|k}
\geq\tfrac14\sum_k\ell_k,
\]
since the least positive system energy gap is $1$.
The ground-to-ground block $T_k$ of $U$ is a contraction
with $\|T_k\|_{2,g}^2=1-\ell_k$.
Its singular value decomposition gives a unitary $V_k$ with
$\|T_k-V_k\|_{2,g}^2\leq\ell_k$: for each singular value $s\in[0,1]$,
$(1-s)^2\leq1-s^2$.

Purify the initial bath as $g^{-1/2}\sum_{a\in G}\ket a\ket a_R$.
Let $w_k$ be the resulting output vector for system input $\ket k$,
and $z_k$ the vector obtained by applying $\ket k\otimes V_k$ instead.
The part of $w_k$ outside $\ket k\otimes G\otimes R$ has squared norm
$\ell_k$, so $\|w_k-z_k\|^2\leq2\ell_k$.
For the uniform pure input, Cauchy--Schwarz gives
\[
\left\|\tfrac12\sum_k(w_k-z_k)\right\|^2
\leq2\sum_k\ell_k\leq8L.
\]
Both summed vectors are unit vectors. Their reduced states are $\tau$
and $D/4$, respectively, so trace-norm contraction under partial trace
and the pure-state trace-distance formula imply
$\|\tau-D/4\|_1\leq2\sqrt{8L}$.
Put $\varepsilon=\|\tau-\sigma\|_1$.
Since $\rho$ and $\sigma$ have the same energy and
$\|H_S-\tfrac72I\|_\infty=7/2$, we have $L\leq7\varepsilon/2$.
Consequently
\[
\max_{j,k}|D_{jk}-C_{jk}|
\leq4\varepsilon+16\sqrt{7\varepsilon}.
\]
For $\varepsilon\leq10^{-8}$ the right side is strictly less than
$1/144$, contradicting Lemma~\ref{lem:gap}.
\end{proof}

\begin{corollary}\label{cor:tail}
For a finite Gibbs bath with excited-state probability $p$, every
energy-conserving output $\tau$ on the input \eqref{eq:data} satisfies
$\|\tau-\sigma\|_1>10^{-8}-2p$.
An exact realization therefore requires $p>5\cdot10^{-9}$.
\end{corollary}
\begin{proof}
The Gibbs state is $(1-p)P_0/g+p\omega_{\rm exc}$, where
$\omega_{\rm exc}$ is a density matrix supported outside the ground space
when $p>0$. Its trace-norm distance from $P_0/g$ is $2p$.
Trace-norm contraction gives $\|\tau-\tau_0\|_1\leq2p$ for the output
$\tau_0$ of the same unitary with bath $P_0/g$.
The triangle inequality and Corollary~\ref{cor:robust} give the result;
the case $p=0$ is Corollary~\ref{cor:robust} itself.
\end{proof}

\section{Sibling statements and extensions}
\begin{corollary}\label{cor:family}
Replace $3/5,4/5$ in $B$ by any real $a,b>0$ with $a^2+b^2=1$, replace
$i$ in its last entry by $e^{i\theta}$ with $\sin\theta\ne0$, and replace
$\ket+$ by any normalized $\sum_j\alpha_j\ket j$ with all $\alpha_j\ne0$.
For every strictly ordered four-level spectrum, the resulting diagonal
cooling channel sends this input to a state that no exact finite
ground-space thermal operation reaches.
\end{corollary}
\begin{proof}
The output entries are $\alpha_j\overline{\alpha_k}C_{jk}$, and its
populations are unchanged. Lemma~\ref{lem:blocks} again forces the Gram
representation. The residual identities give
$V_2=aV_0+bV_1$ and $V_3=aV_0+be^{i\theta}V_1$.
Unitarity gives $X+X^\dagger=0$ and
$e^{-i\theta}X+e^{i\theta}X^\dagger=0$.
Subtracting yields $-2i\sin\theta\,X=0$, the same contradiction.
\end{proof}

\begin{corollary}\label{cor:dimensions}
For every finite system with distinct ordered energies and every pure
input with at least four nonzero energy amplitudes, there is a diagonal
cooling channel whose output lies outside the closure of outputs of
\eqref{eq:thermal}.
\end{corollary}
\begin{proof}
Choose four occupied levels and assign them the rows of $B$; assign
$(1,0)$ to every other level. The two columns define complete diagonal
Kraus operators. Every population is unchanged. Suppose thermal outputs
converged in trace norm to this target. Their energy losses would tend to zero.
The nonnegative energy-loss sum in Lemma~\ref{lem:blocks} then implies
$\ell_k=1-P_{k|k}\to0$ at every occupied level, since its input probability
and every downward energy gap are positive. Complete each ground block to a
unitary as in Corollary~\ref{cor:robust}. The corresponding purified output
vectors differ in norm by at most
$\sum_{k:\alpha_k\ne0}|\alpha_k|\sqrt{2\ell_k}\to0$.
Their reduced states therefore have the same limit. Dividing the coherences
on the four selected levels by their fixed nonzero input amplitudes shows
that the four unitaries' normalized trace Gram matrices converge to $C$,
contradicting Lemma~\ref{lem:gap}.
\end{proof}

\begin{corollary}\label{cor:vectors}
Conjecture 2 of arXiv:1409.7740v1 and v2 (Supplementary Sec.\ II F,
p.\ S11) is false already for three vectors in $\mathbb C^3$.
Property S2.6 in version 3, with its rank-dimensional ground-space
vectors, fails for the rank-two matrix $C$.
\end{corollary}
\begin{proof}
Use the three-vector projection mechanism of
\cite[Proposition 3.2 and Example 3.4]{MR}, also used for exact thermal
separation in \cite[Supplement, Eqs.\ (4)--(10)]{MA}.
Take $v_0=e_0$, $v_1=e_1$, $v_2=(-ie_0+e_1)/\sqrt2$ in $\mathbb C^3$.
If three unitaries $U_t$ made $\{U_tv_j:0\leq t<3\}$ orthonormal for
each $j$, their column matrices $M_j$ would be unitary.
Then $V_j=M_j^\dagger$ would satisfy
$\Tr(V_jV_k^\dagger)/3=v_j^\dagger v_k$.
The zero Gram residual gives $V_2=(iV_0+V_1)/\sqrt2$.
After a common right gauge set $V_0=I$. Unitarity of $V_2$ forces
$V_1=V_1^\dagger$, while $v_1^\dagger v_0=0$ gives $\Tr V_1=0$.
Thus $(I+V_1)/2$ is a projection of rank $3/2$, impossible.
For Property S2.6 the same construction uses the prescribed two
orthogonal unit vectors as the columns of a $2\times2$ matrix $M_j$;
again $V_j=M_j^\dagger$ would represent $C$ as unitary moments.
\end{proof}

\begin{corollary}\label{cor:mixtures}
The channel \eqref{eq:kraus} is outside the closed convex hull of
optimally coherent cooling channels.
\end{corollary}
\begin{proof}
Corollary S2.10 of \cite{NG} puts every finite convex mixture of those
channels in the closure of finite ground-space thermal channels.
If such mixtures converged to \eqref{eq:kraus}, their outputs on $\rho$
would converge to $\sigma$, contradicting Corollary~\ref{cor:robust}.
This uses, and does not refute, the source's approximation result.
\end{proof}

\begin{corollary}\label{cor:threshold}
Under distinct-level and distinct-positive-gap assumptions, the minimum
system dimension admitting a cooling transition outside the closure of
finite ground-space thermal outputs is four.
\end{corollary}
\begin{proof}
The four-level witness is Corollary~\ref{cor:robust}. Qubits satisfy exact
equivalence by \cite[Corollary S2.7]{NG}. For qutrits we recall the closure
proof of \cite[Supplement, pp.\ 1--2]{MA}. Write a cooling channel's
coherence matrix as $Q$, with $q_k=Q_{kk}=P_{k|k}$. Delete the zero
rows and columns; positivity makes them identically zero. The matrix
$R_{jk}=Q_{jk}/\sqrt{q_jq_k}$ on the remaining indices is a correlation
matrix of size at most three. By \cite[Corollary 2.8]{DJ},
$R=\sum_\ell p_\ell z^{(\ell)}z^{(\ell)\dagger}$, where every entry
of each $z^{(\ell)}$ has modulus one. Reinsert the zero rows using
$\sqrt{q_k}=0$. Each term defines an optimally coherent cooling channel
with diagonal Kraus entries $\sqrt{q_k}z_k^{(\ell)}$ and the same lowering
probabilities $P_{j|k}$. The original channel is their convex mixture.
Corollary S2.10 of \cite{NG} puts it in the diamond-norm closure of
finite ground-space thermal channels, so every qutrit output is
approximable. The one-level case is immediate.
\end{proof}

For real correlation matrices
in any dimension, the Clifford matrices in \cite[Theorem 3.1]{DJ} give
an exact finite unitary-moment representation. The controlled unitary
$\sum_j\ket j\bra j\otimes V_j$, with a fully degenerate ancilla,
then realizes the diagonal cooling channel exactly.

\begin{center}
\small
\begin{tabular}{@{}p{.34\linewidth}p{.60\linewidth}@{}}
\toprule
Source item & Coverage \\
\midrule
\cite{NG}, Sec.\ S2 F, Conjecture 1, p.\ S11 & Reproved; prior exact separation Maestro (2026). Refuted (explicit counterexample), strengthened here to closure separation; Theorem~\ref{thm:counterexample} and Corollary~\ref{cor:robust}.\\[3pt]
\cite{NG1,NG2}, Sec.\ II F, Conjecture 2, p.\ S11 & Refuted (explicit three-vector counterexample); Corollary~\ref{cor:vectors}, using the prior projection mechanism of \cite{MR,MA}.\\[3pt]
\cite{NG}, Property S2.6, p.\ S12 & Not asserted universally in the source; fails for the rank-two $C$ (Corollary~\ref{cor:vectors}); sufficiency retained.\\[3pt]
\cite{NG}, Discussion, p.\ 6 & Proved: cooling outputs need not be approximable by finite ground-space thermal outputs; Corollary~\ref{cor:robust}.\\[3pt]
\cite{NG}, Corollary S2.10, p.\ S14 & Prior proof cited; approximation of mixtures retained. Their universality is refuted by Corollary~\ref{cor:mixtures}.\\[3pt]
\cite{NG}, Corollaries S2.7--S2.9 & Source's qubit and rank-one results are retained as prior results; not reproved here.\\[3pt]
\cite{NG}, Theorem S2.4; Sec.\ S3 & Prior proof cited; the cooling criterion and Gibbs-preserving results remain intact.\\[3pt]
\cite{MA}, Letter p.\ 1; Supplement pp.\ 1--2 & Prior proof cited: exact qutrit separation (irrational weight), and qutrit closure. The closure proof is recalled in Corollary~\ref{cor:threshold}; together with Corollary~\ref{cor:robust}, it gives the four-level threshold.\\[3pt]
\cite{HM}, Remark 4.8(2); \cite{DJ}, Corollary 2.11 & Reproved; prior obstruction K\"ummerer (1986, recorded in \cite{HM}), then Dykema--Juschenko (2009) and Haagerup--Musat (2010/2011). The rational variant has a self-contained proof.\\[3pt]
\cite{DJ}, Corollary 2.8; Theorem 3.1 & Prior proofs cited: the correlation decomposition used for qutrit closure, and real-correlation exact realization.\\[3pt]
\cite{DJ}, Example 4.3 & Not covered: the exact depolarizing boundary for their different square-root example.\\[3pt]
Other models & Proved: small Gibbs-tail bound, Corollary~\ref{cor:tail}. Not covered: arbitrary positive-temperature conversion \cite{DHD,CS}, infinite ancillas and catalysts.\\
\bottomrule
\end{tabular}
\end{center}

\paragraph{Mathematical verification.}
The Lean declaration \texttt{result} in the
\href{https://github.com/the-omega-institute/trureturing/blob/0eae8e045cb2ba958e2097dbac034f62f9c1dd7a/D5/S3/QuantumChannels/CoolingMapThermalOperationRefutation.lean}{pinned Lean source}
has statement $\neg\,\texttt{claim}$. Its definitions encode diagonal or
single-lowering Kraus operators and all finite ancillas in the normalized
full ground space. Its witness is exactly \eqref{eq:data}, with the two
operators \eqref{eq:kraus}. Its exact scope is Theorem~\ref{thm:counterexample},
with energies $0,1,3,7$ and the uniform pure input. The uniform Gram gap,
trace-norm separation, parameter and dimension extensions, and vector and
mixture consequences, including the Gibbs-tail bound, have written proofs
here, not additional conclusions
of \texttt{result}. The accompanying Python
checker uses rational complex arithmetic for the concrete Gram matrix,
Kraus completeness and output, density-matrix minors, gap distinctness,
two zero residual norms, nonzero obstruction determinant, and the rational
constants in Lemma~\ref{lem:gap} and Corollary~\ref{cor:robust}.

\begin{thebibliography}{12}
\bibitem{NG} V. Narasimhachar and G. Gour,
\emph{Low-temperature thermodynamics with quantum coherence},
Nature Communications \textbf{6}, 7689 (2015).
\href{https://doi.org/10.1038/ncomms8689}{doi:10.1038/ncomms8689}.
\href{https://arxiv.org/abs/1409.7740v3}{arXiv:1409.7740v3}, Supplementary Sec.\ S2 F.
\bibitem{NG1} V. Narasimhachar and G. Gour,
\emph{Low-temperature thermodynamics with quantum coherence},
\href{https://arxiv.org/abs/1409.7740v1}{arXiv:1409.7740v1}
(26 September 2014), Supplementary Sec.\ II F, Conjecture 2, p.\ S11.
\bibitem{NG2} V. Narasimhachar and G. Gour, same title,
\href{https://arxiv.org/abs/1409.7740v2}{arXiv:1409.7740v2}
(2014), Supplementary Sec.\ II F, Conjecture 2, p.\ S11.
\bibitem{DJ} K. Dykema and K. Juschenko,
\emph{Matrices of unitary moments}, Mathematica Scandinavica \textbf{109},
225--239 (2011); preprint (2009).
\href{https://doi.org/10.7146/math.scand.a-15186}{doi:10.7146/math.scand.a-15186}.
\href{https://arxiv.org/abs/0901.0288}{arXiv:0901.0288}, Proposition 2.10 and Corollary 2.11.
\bibitem{HM} U. Haagerup and M. Musat,
\emph{Factorization and dilation problems for completely positive maps on
von Neumann algebras}, Communications in Mathematical Physics
\textbf{303}, 555--594 (2011); preprint (2010).
\href{https://doi.org/10.1007/s00220-011-1216-y}{doi:10.1007/s00220-011-1216-y}.
\href{https://arxiv.org/abs/1009.0778}{arXiv:1009.0778}.
\bibitem{MR} M. Musat and M. R\o rdam,
\emph{Non-closure of quantum correlation matrices and factorizable channels
that require infinite dimensional ancilla}, Communications in Mathematical
Physics \textbf{375}, 1761--1776 (2020; online 2019); preprint (2018),
with an appendix by N.\ Ozawa.
\href{https://doi.org/10.1007/s00220-019-03449-w}{doi:10.1007/s00220-019-03449-w}.
\href{https://arxiv.org/abs/1806.10242v4}{arXiv:1806.10242v4},
Proposition 3.2 and Example 3.4.
\bibitem{MA} Maestro v0.1, \emph{Minimum Bath Dimension for Qutrit Cooling Maps},
27 August 2026, v1.0, Letter and Supplemental Material.
\href{https://pudimphysics.org/research/RO-2026-0004/releases/v1.0/combined.pdf?download=1}{Full scientific text},
Letter p.\ 1 and Supplement pp.\ 1--2.
\bibitem{MSE} V.\ Narasimhachar,
\emph{Question about existence of unitary matrices with certain properties},
Mathematics Stack Exchange, question 896675 (13 August 2014).
\href{https://math.stackexchange.com/questions/896675/question-about-existence-of-unitary-matrices-with-certain-properties}{Original question}.
\bibitem{DHD} Y.\ Ding, F.\ Ding and X.\ Hu,
\emph{Exploring the gap between thermal operations and enhanced thermal operations},
Phys.\ Rev.\ A \textbf{103}, 052214 (2021).
\href{https://doi.org/10.1103/PhysRevA.103.052214}{doi:10.1103/PhysRevA.103.052214}.
\href{https://arxiv.org/abs/2101.05690v2}{arXiv:2101.05690v2}.
\bibitem{LSBNW} S.\ H.\ Lie, J.\ Son, P.\ Boes, N.\ H.\ Y.\ Ng and H.\ Wilming,
\emph{Thermal operations from informational equilibrium} (2025).
\href{https://arxiv.org/abs/2507.16637v1}{arXiv:2507.16637v1}, Appendix C.2.
\bibitem{MUS} M.\ Musat,
\emph{The Connes--Kirchberg Problem and infinite-dimensional phenomena in
quantum information theory} (2025).
\href{https://arxiv.org/abs/2512.00432v1}{arXiv:2512.00432v1}, Proposition 3.6.
\bibitem{CS} P.\ \'Cwikli\'nski and M.\ Studzi\'nski,
\emph{Limitations on joint coherence preservation in quantum thermodynamics},
version 3 (26 October 2026).
\href{https://arxiv.org/abs/2608.11872v3}{arXiv:2608.11872v3},
Sec.\ III B and Appendix E.5.
\end{thebibliography}
\end{document}
